\documentclass[a4paper,11pt]{jsarticle}

\usepackage[english]{babel}
\usepackage{amssymb, amsmath, amsthm}
\usepackage{bm}
\usepackage[dvipdfmx]{graphicx}
\usepackage{authblk}
\usepackage{indentfirst}
\usepackage{mathtools}
\usepackage{enumitem}

\theoremstyle{definition}
\newtheorem{dfn}{Definition}[section]
\newtheorem{prop}[dfn]{Proposition}
\newtheorem{lem}[dfn]{Lemma}
\newtheorem{thm}[dfn]{Theorem}
\newtheorem{cor}[dfn]{Corollary}
\newtheorem{rem}[dfn]{Remark}
\newtheorem{fact}[dfn]{Fact}

\renewcommand{\refname}{References}

% (Useful) Sets
\newcommand{\rnum}{\mathbb{R}}
\newcommand{\NaturalNumberSet}{\mathbb{N}}
\newcommand{\RealNumberSet}{\mathbb{R}}
\newcommand{\NDemenstionalRealEuclideanSpace}{\mathbb{R}^n}

% (Useful) Texts
\newcommand{\SuchThat}{\quad\text{s.t.}\quad}
\newcommand{\With}{\quad\text{with}\quad}

\newcommand{\Interior}[1]{\mbox{\rm int}\mbox{$\,#1$}}
\newcommand{\Closure}[1]{\mbox{\rm cl}\mbox{$\,#1$}}
\newcommand{\Boundary}[1]{\mbox{\rm bd}\mbox{$\,#1$}}
\newcommand{\Complement}[1]{\mbox{$\,#1^{c}$}}
\newcommand{\OrderingLevelSets}[3]{\text{\rm lev\:$({#1}, {#2}, {#3})$}} % lev


\newcommand{\Painleve}{Painlev\'e}
\newcommand{\InnerProduct}[2]{\left\langle {#1},{#2}\right\rangle} % <x,y>
\newcommand{\Norm}[1]{\left\lVert {#1} \right\rVert} % <x,y>
\newcommand{\SetValudMapping}{\it \Gamma}
\newcommand{\pow}[1]{\mathcal{P}_{0}(#1)}

% Set form e.g. {x | ...}
% #1: element
% #2: conditions
\newcommand{\SetForm}[2]{
  \{{#1}\:|\:{#2}\}
}



\SetLabelAlign{Center}{\hfil#1\hfil}
\SetLabelAlign{CenterWithParen}{\hfil(\makebox[1.0em]{#1})\hfil}

\newcommand{\defarrow}{\;\mathrel{\overset{\mathrm{def}}{\Longleftrightarrow}}\;}
\newcommand{\bleq}[1]{\leq_{#1}\!}
\newcommand{\setrel}[2]{\preccurlyeq^{(#1)}_{#2}\!}
\newcommand{\notsetrel}[2]{{\not\preccurlyeq}^{(#1)}_{#2}\!}
\newcommand{\st}{\ \mathrm{s.t.}\ }


\title{%
Scalarization methods with a characterizing function of \\
an ordering cone in set optimization \\
  \large (集合最適化における順序錐による特徴づけ関数を用いたスカラー化
  )}
\author{
  Ryota Iwamoto$^{a}$ and Tamaki Tanaka$^{b}$
  }
\affil{
  $^{a}$Graduate school of Science and Technology, Niigata University, \\
  $^{b}$Faculty of Science, Niigata University \\
}
\date{\vspace{-5ex}}

\begin{document}
\maketitle
\section{\rm Introduction}
In 2024, Dechboon and Tanaka examined whether composite functions of a set-valued
map and several scalarization functions originated from Gerstewitz's (Tammer's) sublinear
scalarizing functional (see \cite{Tammer2003})
are lower (upper) semicontinuous.

Last year, in this workshop, we talked about ``Set-Valued Fan-Takahashi Inequalities
with Set-Relations Based on Scalarization Methods'' (see \cite{IwamotoTanaka2024}).

This paper, containing these results, explains relations between
a characterizing function of an ordering cone and scalarization functions which lead to
set-valued Fan-Takahashi inequalities.

\section{\rm Preliminaries}

Let $X$ be a topological space and $Y$ a real topological vector space
with the zero vector $\theta_Y$.
Define that $\pow{Y}$ is the set of all nonempty subsets of $Y$.
The sets of neighborhoods of $x \in X$ and
$y \in Y$ is denoted by $\mathcal{N}_X (x)$ and $\mathcal{N}_Y (y)$, respectively.

For $A \in \pow{Y}$, the interior, the closure, the boundary, and the complement of $A$ are denoted by
$\Interior{A}$, $\Closure{A}$, $\Boundary{A}$, and $\Complement{A}$, respectively.
For given $A,B\in\pow{Y}$ and $t\in\rnum$,
the algebraic sum $A+B$ and the scalar multiplication $tA$ are denoted as follows:
$A+B:= \{a+b \mid a\in{A},b\in{B} \}, \quad tA:= \{ta\mid a\in{A}\}$.
In particular, we denote $A+\{y\}$ by $A+y$ and
$(-1)A$ by $-A$ for $A\in\pow{Y}$ and $y\in Y$.
Let us recall that $A$ is said to be $C$-bounded if for each neighborhood $U$ of $\theta_Y$
there exists $t > 0$ such that $A \subset tU + C$.

We define set-relations on $\pow{Y}$ as follows;

\begin{dfn}[set-relations, \cite{KuroiwaTanakaHa1997}]\label{SetRelations}
  Let $C$ be a convex cone in $Y$.
  The six types of set-relations are defined as follows:
  \begin{enumerate}
    %
    \item $A\setrel{1}{C}B \defarrow
            \forall{a\in{A}}, \forall{b\in{B}}, a\bleq{C}b
            \\
            \iff
            A\subset{\left(\bigcap_{b\in{B}}(b-C)\right)}
            \iff
            B\subset{\left(\bigcap_{a\in{A}}(a+C)\right)};$
          %
    \item $A\setrel{2L}{C}B \defarrow
            \exists{a\in{A}}\st\forall{b\in{B}}, a\bleq{C}b
            \iff
            A\cap{\left(\bigcap_{b\in{B}}(b-C)\right)}\ne\emptyset$;
          %
    \item $A\setrel{2U}{C}B \defarrow
            \exists{b\in{B}}\st\forall{a\in{A}}, a\bleq{C}b
            \iff
            \left(\bigcap_{a\in{A}}(a+C)\right)\cap{B}\ne\emptyset$;
          %
    \item $A\setrel{3L}{C}B \defarrow
            \forall{b\in{B}}, \exists{a\in{A}}\st a\bleq{C}b
            \iff B\subset{A+C}$;
          %
    \item $A\setrel{3U}{C} B \defarrow
            \forall{a\in{A}}, \exists{b\in{B}}\st a\bleq{C}b
            \iff A\subset{B-C}$;
          %
    \item $A\setrel{4}{C}B \defarrow
            \exists{a\in{A}}, \exists{b\in{B}}\st a\bleq{C}b
            \\
            \iff A\cap{(B-C)}\ne\emptyset
            \iff (A+C)\cap{B}\ne\emptyset$.
  \end{enumerate}
\end{dfn}

For set-valued maps and scalarization functions, we introduce the following concepts of continuity and semicontinuity depending on
binary relations and cones, which are works of Dechboon and Tanaka \cite{DechboonTanaka2024}.
We can obtain the result of upper semicontinuity for composite functions similarly.

\begin{dfn}{(\cite{DechboonTanaka2024})}
  Let $F \colon X \to \pow{Y}$, $x_0 \in X$, $\preccurlyeq$ a binary relation on $\pow{Y}$
  and $C \subset Y$ a convex cone. We say that $F$ is $(\preccurlyeq, C)$-continuous at $x_0$ if
  \begin{equation}
    \forall W \subset Y, W\:\text{open},W \preccurlyeq F(x_0), \exists V \in \mathcal{N}_{X}(x_0) \st W + C \preccurlyeq F(x), \forall x \in V. \notag
  \end{equation}
\end{dfn}

Especially, $(\preccurlyeq_1, C)$-continuity and $(\preccurlyeq_2, C)$-continuity coincide with classical “$C$-lower
continuity” and “$C$-upper continuity” for set-valued maps, respectively.

\begin{dfn}{(\cite{DechboonTanaka2024})}
  Let $\varphi \colon \pow{Y} \to \RealNumberSet \cup \{\pm \infty\}$, $A_0 \in \pow{Y}$,
  $\preccurlyeq$ a binary relation on $\pow{Y}$, and $C$ a convex cone in $Y$ with $C \ne Y$. Then,
  we say that $\varphi$ is $(\preccurlyeq, C)$-lower semicontinuous at $A_0$ if
  \begin{equation}
    \forall r < \varphi (A_0), \exists W \in \pow{Y}, W\:\text{open}, \st W \preccurlyeq A_0 \:\text{and}\:
    r > \varphi (A), \forall A \in U(W + C, \preccurlyeq); \notag
  \end{equation}
  where $U(V,\preccurlyeq) \coloneqq \SetForm{A \in \pow{Y}}{V \preccurlyeq A}$.
\end{dfn}

\begin{thm}{(\cite{DechboonTanaka2024})}
  Let $F \colon X \to \pow{Y}$, $\varphi \colon \pow{Y} \to \RealNumberSet \cup \{\pm \infty\}$, $x_0 \in X$,
  $\preccurlyeq$ a binary relation on $\pow{Y}$, and $C$ a convex cone. If $F$ is $(\preccurlyeq, C)$-continuous at $x_0$
  and $\varphi$ is $(\preccurlyeq, C)$-lower semicontinuous at $F(x_0)$, then $(\varphi \circ F)$ is lower semicontinuous at $x_0$.
\end{thm}

\section{\rm Set-Valued Fan-Takahashi Inequalities}

Let $\varphi: \pow{Y} \to \rnum \cup \{\pm \infty\}$, $\preccurlyeq$ a binary relation
on $\pow{Y}$, and $C' \subset Y$  a convex cone.
To generalize four types of set-valued minimax inequalities \cite{KuwanoTanakaYamada2010},
we provide a new class of scalarization functions that satisfy;
\begin{enumerate}
  \item $\varphi$ is $(\preccurlyeq, C')$-lower semicontinuous,
  \item $\varphi$ is quasi concave,
  \item $\varphi(\{\theta_{Y}\}) = 0$.
\end{enumerate}
In addition, we give necessary conditions between inequalities and set-relations as follows;
\begin{enumerate}[label=(B\arabic*)]
  \item $\varphi$ is $(\setrel{j}{\Interior{C}})$-monotone,
  \item $\varphi (A) > 0 \Rightarrow  \{\theta_Y\} \setrel{j}{\Interior{C}} A \quad\text{for any}\:A\in \pow{Y}$,
\end{enumerate}
where $C$ is a convex cone with $\Interior{C} \ne \emptyset$ in $Y$ and $j= 1, 2L, 2U, 3U, 3L, 4$.
If $\varphi$ satisfies conditions (i)--(iii), (B1), and (B2),
we write the notation as $\varphi \in \Phi(\setrel{j}{\Interior{C}}\:, \preccurlyeq, C')$.

\begin{thm}\label{NewSetValuedFanTakahashiMinimaxInequalitiesType2}
  Let $X$ be a nonempty compact convex subset of a Hausdorff topological vector space,
  $Y$ a real topological vector space, $\preccurlyeq$ a binary relation on $\pow{Y}$,
  $C$ a convex cone in $Y$, $C'$ a convex cone in $Y$,
  $\varphi\colon \pow{Y} \to \rnum \cup \{\pm \infty\}$,
  and $F\colon X \times X \to \pow{Y}$.
  For the scalarization function $\varphi \in \Phi(\setrel{j}{\Interior{C}}\:, \preccurlyeq, C')$
  satisfying Assumption (A2) where $j= 3U, 3L$,
  if $F$ satisfies the following conditions:
  \begin{enumerate}
    \item $(\varphi \circ F)(x,y) \in \rnum$ for all $x,y \in X$,
    \item for each fixed $y \in X$, $F(\cdot,y)$ is $(\preccurlyeq, C')$-continuous,
    \item for each fixed $x \in X$, $F(x,\cdot)$ is type $(j)$ naturally quasi $C$-concave,
    \item for all $x \in X$, $\{\theta_Y\} \notsetrel{j}{\Interior{C}}\: F(x,x) $,
  \end{enumerate}
  then there exists $\bar{x} \in X$ such that $\{\theta_Y\} \notsetrel{j}{\Interior{C}}\: F(\bar{x},y)$
  for all $y \in X$.
\end{thm}

\section{\rm Characterizing Functions of an Ordering Cone and Some Related Results}

In this section, we describe the definition of characterizing functions.

\begin{dfn}{(\cite{HuqTanaka2025})}
  Let $C$ be a convex cone in $Y$ and $f: Y \to \rnum \cup \{\pm \infty\}$.
  We say that $f$ is a characterizing function of the ordering cone $-C$ if
  \begin{equation}
    f(y) \leq 0 \Leftrightarrow y \in -C \quad \text{for all}\: y \in Y. \notag
  \end{equation}
\end{dfn}

We define a class of scalarization functions for set-valued
Fan-Takahashi inequalities, characterized by
monotonicity, continuity, and ordering with the origin.

\begin{dfn}\label{Scalarization3ConditionsDefinition}
  Let $C$ be a convex cone in $Y$ and
  $\setrel{j}{C} (j= 1, 2U, 2L, 3U, 3L, 4)$ set-relations in Definition~\ref{SetRelations},
  For a scalarization function $\varphi: \pow{Y} \to \rnum \cup \{\pm \infty\}$,
  we consider a class of scalarization functions satisfying the
  following three conditions is denoted by $\Phi_0(\setrel{j}{C})$;
  \begin{enumerate}
    \item $\varphi$ is $(\setrel{j}{C})$-monotone,
    \item $\varphi (A) \leq 0 \Rightarrow  A \setrel{j}{C} \{\theta_Y\}\quad\text{for any}\:A\in \pow{Y}$,
    \item $\varphi(\{\theta_{Y}\}) \leq 0$.
  \end{enumerate}
\end{dfn}

\begin{lem}
  Let $C$ be a convex cone in $Y$,
  $\setrel{j}{C} (j= 1, 2U, 2L, 3U, 3L, 4)$ set-relations in Definition~\ref{SetRelations},
  and a scalarization function $\varphi \in \Phi_0(\setrel{j}{C})$.
  Then, if we define $f(y) = \varphi (\{y\})$ for any $y \in Y$,
  $f$ is a characterizing function of $-C$.
\end{lem}

\begin{lem}
  Let $C$ be a closed convex cone in $Y$ and $f: Y \to \rnum \cup \{\pm \infty\}$
  a characterizing function of $-C$. Define $\varphi_s : \pow{Y} \to \rnum \cup \{\pm \infty\}$ by
  \begin{equation}
    \varphi_s (A) = \sup_{a \in A} f(a) \quad \text{for all}\: A \in \pow{Y}. \notag
  \end{equation}
  If $f$ satisfies that for all $a, b \in Y$ with $a \bleq{C} b$, $f(a) \leq f(b)$,
  then $\varphi_s$ is an element of $\Phi_0(\setrel{3U}{C})$.
\end{lem}

\section{\rm Conclusion}


We can give an example of interior ordering type of set-valued minimax
inequalities.
By observing the new class of scalarization functions in the beginning of section 3,
we identified several properties related to characterizing functions of an ordering cone (see Definition 4.2).
These properties serve as key elements for establishing a new example for the
set-valued Fan-Takahashi inequalities.

\begin{thebibliography}{9}
  \bibitem{DechboonTanaka2024}
  P.~Dechboon and T.~Tanaka,
  Inheritance properties on cone continuity for set-valued maps via
  scalarization, Minimax Theory Appl. {\bf 9} (2), (2024), 201--224.

  \bibitem{Tammer2003}
  A. G\"{o}phert, H. Riahi, C. Tammer, C. Z\u{a}linescu: Variational Methods in Partially Ordered
  Spaces, CMS Books in Mathematics vol. 17, Springer, New York (2003).

  \bibitem{HuqTanaka2025}
  M.\ Z.\ Huq and T.\ Tanaka, {\em Characterization of set-relations in set optimization via twofold scalar optimization problems},
  Taiwanese J. Math. {\bf29} (4), (2025), 825--849

  \bibitem{IwamotoTanaka2024}
  R.Iwamoto and T.Tanaka,
  Set-valued Fan-Takahashi inequalities via scalarization,
  RIMS Kokyuroku, {\bf 2323} (2025), 155--159.

  \bibitem{KuwanoTanakaYamada2010}
  I.~Kuwano, T.~Tanaka and S.~Yamada,
  Unified scalarization for sets and set-valued Ky Fan minimax inequality,
  J. Nonlinear Convex Anal. {\bf 11} (3), (2010), 513--525.

  \bibitem{KuroiwaTanakaHa1997}
  D.~Kuroiwa, T.~Tanaka and T.~X.~D.~Ha,
  {\em On cone convexity of set-valued maps},
  Nonlinear Anal.
    {\bf 30} (1997), 1487--1496.
\end{thebibliography}
\end{document}
